\section{Sprint 0}

 

\subsection*{Summary}

Sprint 0 was started on 15 February and was used as a preparation sprint to set up the entire project infrastructure, including the project repository in GitHub, the project management board in Jira for the SCRUM process, and communication channels as a WhatsApp group. The main goal was to establish the Scrum process and define the initial system structure before starting implementation in sprint 1. 

Also in this sprint, the team had its first meeting with the supervisor and distributed roles between members. A project backlog with user stories and tasks, epics, and assigned them to the team, and made the first UI concept.

\subsection*{Deliverables}

The main deliverable of Sprint 0 was the project backlog with user stories and tasks, epics, and they were assigned them to the team. Also, the project repository was created in GitHub, and the project management board in Jira for the SCRUM process. The team tasks on sprint 0 were:

\begin{itemize}

    \item Create a draft how the product should look like

    \item  Define three-layer architecture

    \item Define high-level UI structure

    \item Define UML class

    \item Jira board and initial Product Backlog
    
    \item GitHub repository

\end{itemize}

 

\subsection*{Sprint Planning}

During Sprint planning, the team chose to start work on setup tasks and workflow preparation,flow, and left coding for the next sprints. The main focus was on app architecture and Scrum process. 

Also, the team discussed what the project should look like. In the Jira workflow, it was decided to make epics for major project areas, so tasks and User Stories could be grouped more clearly and transparently. Also, distributed roles for the first demo presentation. Evidence of Sprint 0 Planning and Review can be found in Figure~\ref{fig:sprint0-planning-review}.



 

\subsection*{Sprint Review}

During the sprint review, the team presented the project backlog with User Stories and tasks, epics, and how tasks and user stories were assigned to team members. As well, the project repository was created in GitHub, the project management board in Jira for the SCRUM process, and communication lines as a WhatsApp group. The first UI concept of an application was demonstrated. 

The feedback from the Danfoss representative, teacher, and supervisor showed that the team needed to define project requirements more clearly, and the team should finish User-Stories and tasks for all sprints in Backlog.

Based on this feedback, the team decided to improve the backlog structure and use Jire more actively in the following sprints. Evidence of Sprint 0 Planning and Review can be found in Figure~\ref{fig:sprint0-planning-review}.

 

\subsection*{Sprint Retrospective}

What went well:

\begin{itemize}

 \item Basic project infrastructure was successfully created

 \item GitHub, Jira, and communication channels were set up

 \item The first Product Backlog and UI concept were prepared

\end{itemize}



What did not go well:

\begin{itemize}

 \item The project scope was still not fully clear
 
 \item The backlog was not detailed enough  
 
 \item The team was still learning how to document Scrum ceremonies properly

\end{itemize}



Improvement:

\begin{itemize}

 \item Improve meeting documentation and Scrum evidence

 \item Refine the Product Backlog in future sprints

 \item Assign tasks more clearly between team members

\end{itemize}

Sprint 0 Retrospective evidence can be found in Figure~\ref{fig:sprint0-retrospective}.

\subsection*{Reflections and Conflicts}

The main difficulty was the uncertainty about project expectations. The team didn’t have a clear view of how the application should work, and realized. As well, the team was not fully confident in the level of detail for Scrum documentation. This question was solved in a meeting with supervisor Alan and after the first demo feedback.

 

% --------------------------------

\section{Sprint 1}

 

\subsection*{Summary}

Sprint 1 started on 2 March and served as an architecture and requirements sprint. The main goal was to focus on defining the system  requirements and UML diagrams before starting implementation and development work in later sprints. 

During this sprint, the team focused on turning the initial idea and concept into a structured design. To complete this work, the group began creating a Figma UI prototype. First, an AI-generated Figma version of the app UI was created to explore possible design approaches. Afterward, a manually designed version was developed by a team member. Additionally, in Sprint 1, daily stand-up meetings were introduced to track progress and discuss task status, as part of implementing the Scrum process. In addition, the backlog was fully implemented, including all User Stories and tasks. The purpose of Sprint 1 was to start technical preparation for the next implementation.

 

\subsection*{Deliverables}

The main deliverables of Sprint 1 were preparing the technical setup, such as diagrams, functional/non-functional requirements, and a UI prototype for the following app development:

 \begin{itemize}
    \item Define functional and non-functional requirements
    \item Create Diagrams (UML, User case, Sequence, Component)
    \item Create a UI-Prototype in Figma

  \item Group Agreement document was created

\end{itemize}


\subsection*{Sprint Planning}

During Sprint Planning, the team discussed and defined the system requirements and architecture. Tasks were then the distributed between team members,such as defining functional and non-functional requirements, creating the all-important project diagrams, and developing a UI prototype. The selected backlog items were moved from the Product Backlog into the Sprint Backlog. Evidence of Sprint 1 Planning can be found in Figure~\ref{fig:sprint1-planning}.

 

\subsection*{Sprint Review}

During the Sprint Review, the group presented a Figma UI prototype, a document with high-level, functional, and non-functional requirements, a sequence diagram, a UML User case, and component diagrams.


The feedback from teachers and supervisors showed  that the team did good work, but they should improve their presentation skills and start working on technical aspects. The Sprint Review confirmed that the project was ready to move into implementation-focused sprints. Evidence of Sprint 1 Review can be found in Figure~\ref{fig:sprint1-review} and Figure~\ref{fig:sprint1-review-board}.

 

\subsection*{Sprint Retrospective}

What went well:

\begin{itemize}

 \item Functional and non-functional requirements were defined

 \item The first diagrams and UI prototype were created

 \item The team improved the Scrum workflow and daily stand-up meetings

\end{itemize}


What did not go well:

\begin{itemize}
    \item Some User Stories were unclear

    \item The backlog required more refinement and details

    \item The team needed better preparation for sprint presentations

\end{itemize}

Improvement:

\begin{itemize}
    \item Improve backlog refinement in Jira

    \item Prepare presentation materials earlier for future sprint reviews

\end{itemize}

Sprint 1 Retrospective evidence can be found in Figure~\ref{fig:sprint1-retrospective}.

\subsection*{Reflections and Conflicts}

During Sprint 1, there were no serious conflicts or problems. The issue of unclear User Stories was improved through backlog refinement in Jira. But in the next sprint presentation, the team decided to make better preparations for the Sprint 2 Demo to make it more professional.

 

% --------------------------------

\section{Sprint 2}

 

\subsection*{Summary}

Sprint 2 started on 16 March and was the start of program development. The main goal was to focus on Asset Manager, Source Data Manager, Optimizer base, and the first UI implementation for implementing scenario 1.

During this sprint, the team focused on using architectural plans to deliver working code. The first stage was to distribute the main system modules and responsibilities among team members. During the daily Scrum meetings, the team discussed component integration and verified that the modules matched the UML class design. One team member also participated in a representative meeting to discuss project progress and communication between groups. The team also demonstrated the first working results using Avalonia UI.

 

\subsection*{Deliverables}

The main deliverables of Sprint 2 were writing code for the main program components and implementing the draft UI:

  \begin{itemize}
    \item Source Data Manager implementation
    
    \item Optimizer logic for Scenario 1 was implemented
    
    \item  Implemented Asset Manager, ProductionUnit class, and UI for displaying production units

    \item Added support for two seasons in SourceDataManager

    \item Optimizer supports different scenarios based on the selected season

\end{itemize}



\subsection*{Sprint Planning}

During Sprint Planning, the team decided to start work on scenario 1. The team moved selected User Stories from the Product Backlog into the Sprint Backlog and identified a dependency chain between the modules: that AssetManager provides units, SourceDataManager provides heat demand/electricity prices, and Optimizer depends on SourceDataManager and AssetManager. And UI will show the current state. Each team member was responsible for a specific module or task. Evidence of Sprint 2 Planning can be found in Figure~\ref{fig:sprint2-planning}.

 

\subsection*{Sprint Review}

During the Sprint Review, the group presented the implemented Asset Manager and Source Data Manager, demonstrated the foundation of the optimizer logic and how the UI looks, and showed that the system structure was beginning to match the architecture defined in the UML class and component diagrams. Additionally, the maintenance functionality was not yet fully implemented. 

The feedback from teachers and the supervisor confirmed that the code and results are correct, but the group received a recommendation to improve the UI. The Sprint Review confirmed that the project was moving from design into implementation-focused development. Evidence of Sprint 2 Review can be found in Figure~\ref{fig:sprint2-review}.

 

\subsection*{Sprint Retrospective}

What went well:

\begin{itemize}

    \item The code runs well and was correctly pushed to GitHub

    \item Everybody kept up with the schedule 

\end{itemize}



What did not go well:

\begin{itemize}

    \item  Team members did not manage to implement the maintenance part for the project and moved it to the next sprint 

\end{itemize}



Improvement:

\begin{itemize}

    \item Team should make sure to take on only as much as they are sure they will be able to accomplish

    \item Additional tasks should only be added after the planned work is completed

\end{itemize}

{fig:sprint2-retrospective}.

\subsection*{Reflections and Conflicts}

The challenge was to distribute story points and prioritize tasks correctly, and some issues with GitHub branch management appeared in the first days of the sprint, but in the daily Scrum meetings, this problem was solved. 

 

% --------------------------------

\section{Sprint 3}

 

\subsection*{Summary}

Sprint 3 started on 27 March and focused on testing and finalizing Scenario 1. The main goal was to focus on creating unit tests for the code, finalizing the remaining features for Scenario 1, and starting to work on Scenario 2.

During this sprint, the team focused on finishing all code for scenario 1. One of the tasks was to complete all the UI so the program would be ready for the midterm evaluation. In addition to extensive unit testing for all main managers and optimizers, a major part of the unit tests included positive/negative/edge cases. Moreover, work started on Scenario 2 and on creating tests for it. In addition, after the LaTeX Workshop team began work on a Report Setup, Visual Studio Code was chosen as the main environment for writing the report in LaTeX.

 

\subsection*{Deliverables}

The important deliverables of Sprint 3 were writing tests for the main program components and implementing the full UI for Scenario 1:

\begin{itemize}
    \item  tests for GetProductionUnits()
    
    \item  tests for SourceDataManager

    \item  tests for RunScenario1()

    \item  tests for RunScenario2()

    \item  tests for GetUnitsByNetProductionCost()

    \item tests for CalculateNetProductionCost()

    \item  tests for CreateMaintenanceForBoiler()

    \item  Completed maintenance logic

    \item  UI for Scenario 1 becomes fully implemented and usable

    \item Completed report repository setup

    \item  Created presentation for midterm evaluation

\end{itemize}

 

\subsection*{Sprint Planning}

This sprint planning was conducted online. During sprint planning, it was decided to focus on unit testing existing methods and improving the UI for Scenario 1 at the end of the meeting. Each member was assigned a Jira testing story. After a recommendation from the supervisor, report-related tasks were added to the Jira backlog. The sprint scope was expanded after the planned tasks were completed earlier than expected. The sprint was more quality-focused than feature-focused. Evidence of Sprint 3 Planning can be found in Figure~\ref{fig:sprint3-planning}.

 

\subsection*{Sprint Review}

During the Midterm Evaluation, the group presented unit testing progress, a fully implemented UI, and the functionality of Scenario 1.

Feedback from teachers, the supervisor, and the Danfoss representative indicated that the graph visualization and UI presentation should be improved. All tasks planned during Sprint Planning were completed before the midterm evaluation. Evidence of Sprint 3 Review can be found in Figure~\ref{fig:sprint3-review}.

 

\subsection*{Sprint Retrospective}

What went well:

\begin{itemize}

    \item The code runs well

    \item Everybody kept up with the schedule 

    \item Unit tests increased confidence in the system

    \item The UI for Scenario 1 became fully usable

\end{itemize}



What did not go well:

\begin{itemize}

    \item  Minor time management issues
    
    \item  It was the team's first experience with the online Sprint Planning meetings
    
    \item  Scope changed during the sprint

\end{itemize}



Improvement:

\begin{itemize}

 \item Take only a realistic scope of work

 \item Improve time management to reduce delays

\end{itemize}

Sprint 3 Retrospective evidence can be found in Figure~\ref{fig:sprint3-retrospective}.

 

\subsection*{Reflections and Conflicts}

The presentation for the midterm evaluation was somewhat confusing. The team adjusted speaking roles to avoid repeating information during the presentation. The team also forgot to formally close the sprint on time. As a result, the sprint was completed after the midterm. Moreover, the team forgot to send the sprint 3 documentation to itslearning. To solve this problem, the group decided to send the documents for sprint 3 and sprint 4 together with the Sprint 4 submission.

 

% --------------------------------

\section{Sprint 4}

 

\subsection*{Summary}

Sprint 4 started on 20 April and focused on report writing and finalizing Scenario 2. The main goal was to finish the full program so that Scenarios 1 and 2 function correctly and to add more graphs to make the visualization clearer and easier to understand

During this sprint, the team focused on finalizing the remaining program functionality. One of the biggest improvements was adding Heat Demand graphs, CO2 emissions, and net production cost. In this sprint, the program became more prepared for the final demo, and the team continued preparing the program for the final demo.

\subsection*{Deliverables}

The important deliverables of Sprint 4 were the visualization part of the program: 

 \begin{itemize}

    \item  Report writing continued;
    
    \item  Heat Demand graph for Scenario 1

    \item  Heat Demand graph for Scenario 2

    \item  CO2 emissions graph

    \item  Net Production Cost graph

    \item UI restructuring based on midterm feedback

    \item  Tests for CreateMaintenanceForBoiler()

    \item  CO2 emissions integration into the optimizer

    \item  Result visualization improvements

\end{itemize}


 

\subsection*{Sprint Planning}

Sprint 4 planning started on 20 April, after the midterm, in an offline mode. During this sprint, the team distributed UI visualization tasks and checked the status of the report. Also, the midterm feedback was discussed, along with how to implement all programs with features in the sprint timeframe. Moreover, after the supervisor’s recommendation, the team added a Jira backlog task related to Avalonia unit testing. Evidence of Sprint 4 Planning can be found in Figure~\ref{fig:sprint4-planning}.

 

\subsection*{Sprint Review}

During the Sprint 4 demo, the group presented the results of their work, including new graphs. In addition, using the presentation techniques learned during the additional presentation skills session.

Feedback from teachers and the supervisor confirmed that the work was completed correctly, but the graph should be more colorized and user-friendly for workers, and the team received recommendations to add more unit tests and improve UI. Evidence of Sprint 4 Review can be found in Figure~\ref{fig:sprint4-review}.

 

\subsection*{Sprint Retrospective}

What went well:

\begin{itemize}

    \item The team successfully implemented the midterm feedback

    \item graphs were added

    \item report work continued

    \item UI improved

\end{itemize}



What did not go well:

\begin{itemize}

    \item  Minor time management issues due to the Math exam

    \item  Report writing and development overlapped during the sprint

\end{itemize}



Improvement:

\begin{itemize}

    \item Split work more equally

    \item Continue and finalize report writing

\end{itemize}

Sprint 4 Retrospective evidence can be found in Figure~\ref{fig:sprint4-retrospective}.
 

\subsection*{Reflections and Conflicts}

The main challenge of this sprint was a Math exam scheduled for the same time. Because of the overlapping exam period, the team tried to complete all tasks after the exam within a limited amount of time. The team had to balance project work, report writing, and exam preparation at the same time.

% --------------------------------

\section{Sprint 5}

 

\subsection*{Summary}

Sprint 5 started on 1 May and focused on finalizing the project and unit tests,completing reports, and updating diagrams. The main goal was to finish all parts of the project, from code and diagrams to report and meeting logs before presenting the final version of the project

During this sprint, the team focused on improving and polishing the UI of the program, for example, improving button visibility and UI interaction feedback, and finishing report chapters and receiving feedback on them. Moreover, checking and improving diagrams of the project to make it accurately matched the implemented code. Finally, receiving feedback from the supervisor on how to improve the visual part of the app.

\subsection*{Deliverables}

The main deliverables of Sprint 5 were the finalizing report and program: 

\begin{itemize}
    \item  all unit tests implemented
    
    \item report close to the final version

    \item  All managers and optimizer components worked correctly

    \item  graph functionality testing

    \item  manual testing

    \item final demo preparation

\end{itemize}


 

\subsection*{Sprint Planning}

Sprint 5 planning started on 4 May in an offline meeting. During this sprint, the team distributed tasks, making the UI clearer and more user-friendly. The team also decided to add additional unit tests for graphs in Scenario, and these tasks were added to the Sprint Backlog. Each team member was assigned responsibility for checking unit tests related to a specific manager or module and adding positive, negative, and edge cases where possible. Also, the team decided to prepare for the final demo and demonstrate how the program works in different conditions. Evidence of Sprint 5 Planning can be found in Figure~\ref{fig:sprint5-planning}.


 

\subsection*{Sprint Review}

During the final demo, the group presented the final version of the application to Danfoss representatives, teachers, and supervisors. In this presentation, the team demonstrated all tabs of the program and the flexibility of graphs.

Feedback from Danfoss representatives, teachers, and the supervisor confirmed that the programs work absolutely correctly, and the graphs looked good. The supervisor gave some various design improvements, such as bigger fonts in Asset Manager, clearer selected scenario/seasons, allowing multiple maintenance units, or warning if demand is not being used, and improving the Data tab as a button/page. Evidence of Sprint 5 Review can be found in Figure~\ref{fig:sprint5-review}.

 

\subsection*{Sprint Retrospective}

What went well:

\begin{itemize}

    \item All unit tests were implemented;

    \item Main branch was stabilized;

    \item Report writing continued successfully

    \item UI improved

\end{itemize}



What did not go well:

\begin{itemize}

    \item Report still unfinished

    \item Diagrams were outdated/missing until late

    \item Some test mismatches occurred after changes in the codebase

\end{itemize}



Improvement:

\begin{itemize}

    \item Focus on finishing the report

    \item Update diagrams

    \item Ensure tests are updated when code changes

\end{itemize}

Sprint 5 Retrospective evidence can be found in Figure~\ref{fig:sprint5-retrospective}.
 

\subsection*{Reflections and Conflicts}

There was an issue that the code needed to be updated at the last moment before the final demo, but this problem was quickly resolved, and the final demo was completed successfully. The final phase required balancing report writing, UI polishing, diagram corrections, and final testing at the same time.